\subsection{Projective row modules and faithful actor images}

\begin{theorem}[Proved: projective row modules with two minimal columns]
\label{CP2:A-projective-row-two-column-all-states}
Let $M$ have a proper rectangular minimal band $K$ with two columns.
Put $R=\mathbb QM$, $A=R/\mathbb QK$, and let
$V=\ker(Re\to\mathbb Q)$ be the full row-augmentation module,
where $e$ is any band idempotent.  Suppose $V$ is projective over $A$.
Then every finite left-$R_3$, right-$M$ biset gives a collage whose
rational degree-two products vanish.  Neither bounded homology nor
a bound on the number of rows, outside elements, or states is required.
\end{theorem}
\begin{proof}
For a nonempty state set, use the actual cofiber and its natural
augmented graph chains from
Theorem~\ref{CP2:A-two-column-two-outside-graph}.
The proof of that theorem, up to its final coefficient lemma,
uses only the two minimal columns.  In particular it already
allows the full row module $V$.

If the two column-image graphs share a vertex, the resulting
$M$-linear augmentation section makes the actual bridge pullback
injective on $H^2(BM;\mathbb Q)$.  Every second cofiber class then
comes from its actual suspension quotient, so all products vanish.
Otherwise the column-image graphs have disjoint supports.  The
injective projective subcomplex $Ce\otimes eR$ has two-term quotient
$D$, a right $A$-complex, and the same two augmentation shifts and
projective-ideal change of rings give
\[
 H_n(Y;\mathbb Q)=H_{n-2}(D\otimes_A^{\mathbf L}V),\qquad n\ge4.
\]
Since $V$ is projective, the tensor is represented by the two-term
ordinary complex $D\otimes_A V$.  Thus $H_4(Y;\mathbb Q)=0$.
This argument uses no boundedness assumption.
The actual rational ring comparison with the original collage
proves the assertion in both cases.

If the state set is empty, the collage is the disjoint union of
$BM$ and $BR_3$.  The left-projective resolution
$0\to\mathbb QK\to R\to A\to0$ and projectivity of $V$ over $A$
give $\operatorname{pd}_R V\le1$.  The sequence
$0\to V\to Re\to\mathbb Q\to0$ then gives
$\operatorname{pd}_R\mathbb Q\le2$, so $H^4(BM;\mathbb Q)=0$.
The other component is rationally acyclic.  This also proves
the empty-state assertion.
\end{proof}

\begin{corollary}[Proved: transfer from projective actor images]
\label{CP2:A-projective-actor-image-cup-transfer}
Suppose a surjective homomorphism $M\to H$ identifies their proper
rectangular minimal bands bijectively, those bands have two columns,
and the full row-augmentation module of $H$ is projective over
$\mathbb QH/\mathbb QK$.
Then every collage over a finite left-$R_3$, right-$M$ biset has
zero rational degree-two products.
\end{corollary}
\begin{proof}
The actual induced state quotient gives an $H^2$ isomorphism of
collages by Proposition~\ref{CP2:A-band-quotient-collage-H2}.
All target products vanish by the preceding theorem, and their
actual pullbacks exhaust the source second cohomology.
There is no need to transfer bounded homology through the quotient.
\end{proof}

\begin{remark*}[Literal projectivity certificates]
For a three-row band, the full row module has dimension two.
With the outside monoid elements as a basis of $A$, let $L_a$
be the left regular matrices and $V_a$ the literal row-augmentation
matrices.  The canonical surjection
\[
 P:A^2\longrightarrow V,\qquad (a,b)\longmapsto a v_1+b v_2
\]
is certified to split by a rational matrix $X$ satisfying
\[
 PX=I_2,\qquad \operatorname{diag}(L_a,L_a)X=XV_a
 \quad\text{for every outside basis element }a.
\]
These finite identities prove projectivity directly.  They do not
depend on a classification of small algebras or on faithfulness of
the linear representation.  The separate checker reconstructs the
literal algebra, module and surjection from the source tables and
verifies every coefficient using exact rational arithmetic.
\end{remark*}
